\scanpage{110}

\[
T\longrightarrow Y
\]
\uline{et} d'un relèvement en $T\longrightarrow X$, ou ce qui revient au $=$, d'une section de $X_T=X\times_YT$ sur $T$. Ceci posé,

\noindent
\begin{minipage}[t]{0.48\linewidth}
\vspace*{\baselineskip}
\diagram{\begin{tikzcd}
X\arrow[d] & X_T\arrow[l]\arrow[d] & \\
Y & T\arrow[l] & T\arrow[ul, "{\unsure{}}"{description,pos=0.18}]\arrow[l,phantom,"=" description]
\end{tikzcd}}
% The question near the diagonal's lower-right end marks the possible second
% arrowhead there. The upper-left head is certain; the short lower-right return
% may be a tail return or a second head. See the whole-diagram audit and doubt.
\end{minipage}\hfill
\begin{minipage}[t]{0.48\linewidth}
on aura un isom can (sous l'hyp. de ``suite exacte d'homotopie'' \deleted{\unsure{)}} faite)
\end{minipage}

\[
(12)\qquad \mathcal E_{X/Y}^T\simeq\Pi_1(X_T;T)\simeq\mathcal E_{X_T}^T
\]
et on aura
\[
(13)\qquad 1\longrightarrow\mathcal E_{X/Y}^T\longrightarrow\mathcal E_X^T\longrightarrow\mathcal E_Y^T\longrightarrow1
\]
(Cette hypothèse sera tjrs satisfaite si $Y=\operatorname{Spec}K$, $K$ un corps, si $X$ est géom. $0$-connexe sur $K$)

Plus généralement, si on a une suite exacte d'homotopie universelle, mais pas néc. avec $\mathcal E_X^T\longrightarrow\mathcal E_Y^T$ surj,

\margin{109}
